\section{Pipelined Epoch Transitions}
\label{sec:transitions}
\subsection{Lagged committees and local boundaries}
The account committee for epoch $e$ is $K_e=C_{e-2}$. That outgoing committee decides $S_e$; the already known successor $C_{e-1}=K_{e+1}$ verifies and installs it. The state $S_e$ derives $C_e$, whose account service starts in epoch $e+2$. Figure~\ref{fig:committees} separates selection, installation, and account voting. Genesis supplies $S_{-1}=S_G$ and $C_{-2}=C_{-1}=C_G$.

After local duration $D$, a validator may end epoch $e$ only when no older epoch remains closing. It stops old account signing, records all complete eligible old blocks as notarized, freezes its report against the closed $S_{e-1}$, and opens $e+1$. Learning a decided checkpoint also stops any remaining signing for its closed epoch. Catch-up verifies predecessors in order and never reopens old vote records. A validator signs nothing new in an epoch it has left; it re-broadcasts what it signed before, so that certificates of the closing epoch can still complete from released votes.

There is at most one closing epoch per validator. If closure outlasts $D$, the next local boundary is deferred; successor voting continues, subject to eligibility. The two-epoch lag gives future members an opportunity to synchronize, not an upper bound on state transfer. Boundaries are local events: validators reach them at different times, and nothing in the protocol requires a synchronized stop.

\begin{figure}[t]
\centering
\begin{tabularx}{\columnwidth}{@{}lYYY@{}}
\toprule
Account epoch & $e$ & $e+1$ & $e+2$\\
\midrule
Voters & $C_{e-2}$ & $C_{e-1}$ & $C_e$\\
Role for $S_e$ & Decide & Verify and install & Selected by $S_e$\\
\bottomrule
\end{tabularx}
\caption{Committee roles are staggered. Closure of $e$ overlaps account epoch $e+1$; the table imposes no wall-clock boundaries.}
\Description{A three-epoch table maps account voters to the committee that decides the checkpoint, the successor that verifies and installs it, and the future committee selected from its state.}
\label{fig:committees}
\end{figure}

\subsection{Immutable signed reports}
\label{sec:reports}
Every correct outgoing validator $i$ signs one header
\[
R_i=\Sign_i\bigl(\ctx_e,\Root(T_i),\Root(G_i)\bigr),
\]
where the context binds the session, closing epoch, predecessor checkpoint, outgoing committee, and reporter identity. A Byzantine identity may equivocate, but a candidate selects at most one of its headers.

Let $L_i$ be the reporter's ledger at freeze time. Its canonical tagged set is
\[
T_i=\{(h(B),s_i(B)):B\in L_i\},\qquad s_i(B)\in\{\tagR,\tagN,\tagF\}.
\]
An $\tagR$ entry is inherited recovery-protected state. An $\tagN$ entry has represented notarization but no locally observed finality. An $\tagF$ entry is final by inherited checkpoint state or by an explicit account certificate for the block or a descendant selecting it. The strongest status wins: $\tagF>\tagN>\tagR$. Each entry also names the block's account position and parent, so that a selection places every block it names without needing the block body. Let $V_i^{(e)}$ contain the hashes of all blocks for which $i$ released an epoch-$e$ account vote before freezing. The second committed set is
\[
G_i=V_i^{(e)}\setminus\keys(T_i).
\]
The roots commit to ledger status and to voted-but-not-ledger membership, not to blocks or certificate objects. Correct reporters retain the frozen sets, vote records, and admission evidence. A later vote arrival, certificate assembly, or successor finalization can change their live view but cannot change the signed report.

\subsection{Usable reports}
A root does not enumerate its members, so a validator that holds every block body still does not know which tagged entries a reporter froze. Reconstruction is a protocol precondition, and it is what the convergence service of Section~\ref{sec:services} provides: $\Conv_j(R_i)$ yields $(T_i,G_i)$ with integrity C1. RAI adds the semantic check. A report is \emph{usable} at $j$ when both sets are reconstructed, every entry is well formed, and its claims are justified by evidence $j$ holds: an $\tagR$ entry matches predecessor recovery state; an $\tagN$ entry has the relevant valid $\NC$ or an inherited notarization record; an $\tagF$ entry has valid inherited checkpoint finality or an explicit account certificate, possibly assembled in an earlier retained epoch; and a $G_i$ membership is backed by $i$'s own signed first vote. Relaying another validator's vote never counts as reporter support. A report naming a block above the first position without its parent, unless the predecessor finalized that block there, is malformed; such a report is never selected, so one Byzantine report cannot make every candidate unconstructible. Missing evidence is fetched from retained data; a Byzantine report with unavailable or malformed material stays unusable. Closure needs only $N-f$ usable reports, and at least that many reporters are correct.

\subsection{Eligible candidates}
\label{sec:manifest}
A candidate selects exactly $N-f$ usable report headers from distinct members of $K_e$, in canonical identity order, forming $Q$. It commits to an evidence manifest $\mathcal M$ listing content-addressed blocks, report reconstructions, exact supporting vote records, and finite admission and dependency witnesses; the bytes are retained protocol data, and the manifest introduces no trusted object. The candidate value is
\[
v_e=(Q,\mu_e,d_e),\qquad \mu_e=H(\mathcal M),\qquad d_e=H(S_e),
\]
with $S_e=\Build(S_{e-1},Q,\mathcal M)$. RAI's eligibility predicate, which the consensus service enforces through A2, is
\[
\Elig_e(v)\iff
\begin{cases}
Q\text{ has }N-f\text{ usable reports of distinct members of }K_e,\\
H(\mathcal M)=\mu,\quad \Build(S_{e-1},Q,\mathcal M)=S\neq\bot,\quad H(S)=d.
\end{cases}
\]
All contexts, report identities, vote domains, and predecessor references must verify. A correct member evaluates $\Elig_e$ only on reconstructed and validated inputs; receiving a digest establishes neither availability nor validity. The manifest removes an ambiguity between a committed report \emph{claim} and the evidence later chosen to validate it: witness bytes are fixed before voting, so the state digest is a function of fixed semantic inputs rather than of certificate arrival order. A proposal that extends a certified candidate inside the consensus instance inherits $Q$ and both digests; it cannot replace the report selection. The manifest is not a license to omit required report material, nor may it promote an $\tagN$ or $G_i$ entry to final because an incidental certificate is present.

\subsection{Deterministic construction}
\label{sec:build}
Let $\Anc(X)$ include $X$ and all its account ancestors. A candidate block is an inherited block, a selected tagged block, or a block named by a selected $G_i$, with the ancestors and fixed dependencies needed to validate it. Eligibility is checked against the closed predecessor or the explicit overlap evidence retained with its admission; an unexplained speculative path is excluded from new recovery counts. Recovery admissibility checks anchor, ancestry, dependencies, and cross-boundary witnesses, but not the candidate's own certificate, which would make Rule~2 circular.

For a block $B$, let $\mathsf{f}_Q(B)$ mean that a selected $\tagF$ entry explicitly names $B$ and its finality witness verifies, and $\mathsf{n}_Q(B)$ that selected $\tagN$ evidence represents an eligible notarization. The fresh reporter support set is
\[
U_Q(B)=\{i:h(B)\in G_i\text{ and $\mathcal M$ verifies $i$'s first vote for $B$}\}.
\]
Inherited $\tagR$ entries add no fresh support. Finality witnesses may be for descendants, but only the named block and its selected ancestors are final-output seeds.

\para{Final closure.}
Start with inherited finality and all blocks satisfying $\mathsf{f}_Q$. Take account-ancestor closure and recursively include the explicit finalized-send dependencies required by those final blocks; every added send has a verified finality proof fixed by the signed receive. Call the result $F^\star$.

\para{Rule 1: represented notarization.}
At each new account position, after excluding conflicts with $F^\star$, retain the unique eligible represented $\NC$ target and its ancestry as a notarization lock. Same-domain account exclusion rules out two conflicting eligible $\NC$s. The lock constrains continuation, not balances.

\para{Rule 2: reporter first votes.}
If there is no represented eligible $\NC$ at that position and exactly one eligible block has $|U_Q(B)|\geq r$, retain that block and its ancestry as a recovery lock. Multiple threshold candidates cause neither an arbitrary winner nor certificate-free finalization; Theorem~\ref{thm:closure} shows that a formable fast-final block cannot coexist with a threshold rival.

\para{Inherited state and discharge.}
Retain predecessor unresolved survivors unless explicit compatible finality excludes them or a matching-origin $\NC$ proves a recovery obligation impossible. Check lock discharge before permitting a conflicting new branch; never infer it from a certificate's newer epoch number. All evidence summaries are computed before any omission. Let $R_Q$ be the surviving inherited unresolved state, selected lock targets, and their ancestor closure after finality and valid discharge.

\para{Rule 3: checkpoint promotion.}
For each account, starting at the tip of $F^\star$, promote the next block only while it is the unique child retained in $R_Q$ and the manifest verifies a closing-epoch $\NC$ for it; stop at the first fork or block without such an $\NC$. The resulting prefix $P_Q$ is candidate-final until the consensus service decides the candidate. Recovery-only blocks are never promoted. Let $F^+$ be the ancestor- and finalized-send-dependency closure of $F^\star\cup P_Q$; the output ledger marks $F^+$ final and retains $R_Q\setminus F^+$.

\begin{algorithm}[t]
\caption{\texorpdfstring{$\Build(S_{e-1},Q,\mathcal M)$}{BuildState}}
\label{alg:build}
\begin{algorithmic}[1]
\State Verify the predecessor, selected identities, report roots, and the committed manifest.
\State Take every selected $(T_i,G_i)$ from $\Conv$; verify the material for all required memberships and claims.
\State Form the parent-closed candidate set; classify admission and eligibility from fixed evidence.
\State Compute $F^\star$ by explicit-finality, ancestor, and finalized-send dependency closure.
\State Compute represented $\NC$ targets and reporter support sets for every position before pruning.
\State Check inherited locks and matching-origin discharge witnesses; reject incompatible claims.
\State Apply Rules 1--2; retain inherited survivors and the ancestor closure of required lock targets, forming $R_Q$.
\State Apply Rule~3 to obtain $P_Q$; compute $F^+$ from $F^\star\cup P_Q$.
\State Retain $R_Q\setminus F^+$ and replay $F^+\setminus F_{e-1}$ once in canonical dependency order.
\State Return $(L_e,\Sigma_e)$; Rule~3 has effect only if this candidate is decided.
\end{algorithmic}
\end{algorithm}

Replay orders parents before children and finalized sends before dependent receives; among independent ready blocks it chooses the least $(a,v,h(B))$. A promoted receive uses the send-finality witness fixed at admission and cannot depend on the checkpoint under construction. A cyclic or semantically invalid graph is rejected; Lemma~\ref{lem:construction} shows that an all-correct usable selection is never rejected. Algorithm~\ref{alg:build} fixes the order of validation, summarization, lock selection, and replay.

\subsection{Installing a decided checkpoint}
\label{sec:install}
Let a validator's current closed base be $S_b$ with finalized set $F_b$, and let $O_i$ be its certified overlay, so that its effective final set is $J_i=F_b\cup O_i$. On an accepted proof $\pi_e$ for $v_e$, the validator recomputes $S_e$ from the inputs $v_e$ names, checks $H(S_e)=d_e$, and installs it with final set $F_e$:
\[
J'_i=J_i\cup F_e,\qquad O'_i=J'_i\setminus F_e.
\]
The union is checked for account compatibility, final dependencies, and valid lock discharge. The validator adopts $S_e$ as its base and materializes the live view as $\Replay(\Sigma_e,O'_i)$. This rebases a certified set of effects; it does not authorize old sends again, and persistent operation identifiers suppress duplicate application and duplicate external notifications.

The installation invariant has three parts. The base's final set grows monotonically and may include checkpoint-promoted blocks. Every overlay block has an explicit account-finality proof for itself or a selecting descendant, and the union with the base is dependency-closed. The union extends each base-final account prefix and respects every undischarged notarization or recovery obligation; a recovery lock superseded by a matching-domain exclusion witness leaves the effective view without editing the signed checkpoint.

Committee $C_{e-1}$ holds no second election. It verifies the outgoing committee's proof, reconstructs the state, derives $C_e$ from the canonical $\Sigma_e$, and rechecks nonfinal early work. Additional live finality never changes that committee selection. Invalidated provisional work is discarded without clearing its persisted first vote or terminal state; already finalized overlap work remains certified and applied. Locks the checkpoint retains are continued, not reopened: a notarization lock gets an instance in the open epoch, a recovery-only lock once its block holds a valid closing-epoch $\NC$, and a rival at a retained position is refused.
